\documentclass[11pt]{article}
\usepackage[utf8]{inputenc}
\usepackage[margin=1in]{geometry}
\usepackage{amsmath}
\usepackage{graphicx}
\usepackage{booktabs}
\usepackage{hyperref}
\usepackage[round]{natbib}
\usepackage{float}

\title{Draphnet: Learning a Network Connecting Drug Biological Effects to Disease Genetics for Predicting and Interpreting Drug--Phenotype Associations}
\author{Author A$^{1}$, Author B$^{1,2}$, Author C$^{2}$, Author D$^{1}$\\[6pt]
$^{1}$Section of Genetic Medicine, University of Chicago\\
$^{2}$Department of Human Genetics}
\date{}

\begin{document}
\maketitle

\begin{abstract}
Understanding how drugs produce unexpected effects on diverse diseases---both harmful side effects and beneficial repurposing opportunities---requires connecting molecular drug action to disease genetics. Here we present Draphnet, a supervised affinity regression model that integrates in vitro drug bioassay data from ToxCast (1391 endpoints across 429 drugs) with genome-wide disease gene associations from PhenomeXcan (${\sim}$200 UK Biobank phenotypes) to learn an interpretable network explaining how drug molecular effects propagate to disease phenotypes. Draphnet substantially outperforms a baseline similarity model in predicting held-out drug--phenotype associations from SIDER (Jaccard distance improvement of 0.14). Drugs sharing Draphnet-predicted biological effects also share known gene targets from DrugBank ($p < 0.001$, rank sum test), supporting biological plausibility. The learned network identifies disease genes impacted by drug targets---including the experimentally validated fenofibrate effector CETP as a shared downstream target of anti-inflammatories and PPAR$\alpha$-agonists. Drug groupings based on shared disease-gene effects reveal biologically meaningful cliques that transcend traditional pharmacological classification. These findings demonstrate that integrating bioassay data with disease genetics in a supervised framework can both predict and explain drug--phenotype associations.
\end{abstract}

\section{Introduction}

Drugs frequently produce unexpected effects on diseases beyond their intended therapeutic targets \citep{ashburn2004drug, pushpakom2019drug}. These effects---whether adverse (side effects) or beneficial (repurposing opportunities)---reflect the propagation of drug molecular actions through biological networks to disease-relevant processes. However, the mechanistic pathways connecting in vitro drug activity to clinical phenotypes remain poorly understood \citep{keiser2009predicting, oprea2018unexplored}.

Existing computational approaches to predicting drug--phenotype associations include connectivity score methods based on transcriptomic perturbation signatures \citep{lamb2006connectivity, subramanian2017next} and matrix completion approaches \citep{natarajan2014inductive, lotfi2018bayesian}. While effective for prediction, these methods generally lack biological interpretability---they cannot explain why a drug affects a particular phenotype.

Two rich but underexploited resources offer complementary views of drug and disease biology. The EPA ToxCast/Tox21 program \citep{richard2018toxcast} provides in vitro bioassay data across 1391 molecular endpoints for thousands of compounds, characterizing drug molecular effects at unprecedented scale. PhenomeXcan \citep{pividori2020phenomexcan} derives gene--disease associations from UK Biobank GWAS summary statistics \citep{bycroft2018ukbiobank} using S-MultiXcan \citep{barbeira2018smultixcan}, providing a genomic profile for each of ${\sim}$200 common phenotypes. By connecting these two data sources through known drug--phenotype associations from SIDER \citep{kuhn2016sider}, we can learn an interpretable model that maps drug molecular endpoints to disease genes.

\section{Results}

\subsection{Premise Validation: Molecular Similarity Predicts Phenotypic Similarity}

Before model fitting, we validated the core premise that drugs with similar molecular endpoints tend to share phenotypes, and that phenotypes with similar genomic profiles tend to co-occur with the same drugs (Table~\ref{tab:premise}).

\begin{table}[H]
\centering
\caption{Premise validation: molecular and genomic similarity predict phenotypic associations.}
\label{tab:premise}
\begin{tabular}{lccc}
\toprule
\textbf{Comparison} & \textbf{Metric} & \textbf{Result} & \textbf{$p$-value} \\
\midrule
ToxCast similarity vs.\ shared SIDER phenotypes & Enrichment & Positive association & $1 \times 10^{-22}$ \\
PhenomeXcan similarity vs.\ drug co-occurrence & Spearman $\rho$ & $-0.11$ (Jaccard dist.) & $<0.001$ \\
\bottomrule
\end{tabular}
\end{table}

\subsection{Model Architecture and Data Dimensions}

Draphnet uses an affinity regression framework: $\hat{\mathbf{Y}} = \mathbf{D} \cdot \mathbf{W} \cdot \mathbf{P}^\top$, where $\mathbf{D}$ encodes drug similarity from ToxCast, $\mathbf{P}$ encodes phenotype similarity from PhenomeXcan, and $\mathbf{W}$ is the learned interaction matrix connecting drug endpoints to disease genes (Table~\ref{tab:matrices}).

\begin{table}[H]
\centering
\caption{Draphnet data matrices and dimensions.}
\label{tab:matrices}
\begin{tabular}{llccc}
\toprule
\textbf{Matrix} & \textbf{Source} & \textbf{Rows} & \textbf{Columns} & \textbf{Completeness} \\
\midrule
$\mathbf{Y}$ (labels) & SIDER & 429 drugs & ${\sim}$200 phenotypes & Sparse binary \\
$\mathbf{D}$ (drug features) & ToxCast Level 5 & 429 drugs & 1391 endpoints & ${\sim}$7\% dense \\
$\mathbf{D}$ (reduced) & SoftImpute & 429 drugs & $k$ components & Complete \\
$\mathbf{P}$ (phenotype features) & PhenomeXcan & Top genes & ${\sim}$200 phenotypes & Complete \\
$\mathbf{W}$ (learned) & Training & $k$ & $n_\mathrm{genes}$ & Complete \\
\bottomrule
\end{tabular}
\end{table}

ToxCast data were preprocessed using SoftImpute \citep{mazumder2010softimpute} for dimensionality reduction and missing value imputation. PhenomeXcan gene scores were computed as signed z-scores from S-MultiXcan p-values, keeping the top half of most variably associated genes.

\subsection{Prediction Performance}

Draphnet substantially outperformed a baseline model that uses only drug--drug and phenotype--phenotype similarity without the learned interaction matrix (Table~\ref{tab:prediction}).

\begin{table}[H]
\centering
\caption{Prediction performance: Jaccard distance between predicted and held-out drug effects.}
\label{tab:prediction}
\begin{tabular}{lccc}
\toprule
\textbf{Model} & \textbf{Side effects} & \textbf{Indications} & \textbf{Mean} \\
\midrule
Baseline (similarity only) & $0.72 \pm 0.04$ & $0.68 \pm 0.05$ & $0.70$ \\
Draphnet & $\mathbf{0.58 \pm 0.03}$ & $\mathbf{0.54 \pm 0.04}$ & $\mathbf{0.56}$ \\
Improvement & $0.14$ & $0.14$ & $0.14$ \\
\bottomrule
\end{tabular}
\end{table}

\subsection{Drug Target Validation Against DrugBank}

To assess biological plausibility, we tested whether drugs sharing known gene targets (from DrugBank; \citealt{wishart2018drugbank}) have more similar Draphnet projections into disease gene space (Table~\ref{tab:target_val}).

\begin{table}[H]
\centering
\caption{Drug target validation: similarity of Draphnet projections for drugs sharing targets.}
\label{tab:target_val}
\begin{tabular}{lccc}
\toprule
\textbf{Analysis} & \textbf{Real projection} & \textbf{Null (permuted)} & \textbf{$p$-value} \\
\midrule
Mean Spearman $\rho$ (shared targets) & $0.34 \pm 0.08$ & $0.02 \pm 0.05$ & $<0.001$ \\
Fraction of targets with $p < 0.05$ & $0.68$ & $0.05$ & $<0.001$ \\
Calibration (real vs.\ permuted) & Enriched & At chance & $<0.001$ \\
\bottomrule
\end{tabular}
\end{table}

\subsection{Discovery of Downstream Disease Gene Effects}

The learned matrix $\mathbf{W}$ projects each drug into disease gene space, revealing which disease driver genes are predicted downstream of each drug's molecular effects (Table~\ref{tab:discoveries}).

\begin{table}[H]
\centering
\caption{Selected drug--target--disease gene associations discovered by Draphnet.}
\label{tab:discoveries}
\begin{tabular}{llll}
\toprule
\textbf{Drug category} & \textbf{DrugBank target} & \textbf{Disease gene} & \textbf{Evidence} \\
\midrule
Anti-inflammatories & COX-1/2 & CETP & Shared downstream effect \\
PPAR$\alpha$-agonists & PPAR$\alpha$ & CETP & Fenofibrate effector (validated) \\
PPAR$\gamma$-agonists & PPAR$\gamma$ & RP11-54O7.17 & Adjacent to PERM1 \\
\bottomrule
\end{tabular}
\end{table}

The convergence of anti-inflammatories and PPAR$\alpha$-agonists on CETP is biologically meaningful: CETP (cholesteryl ester transfer protein) has been experimentally validated as a downstream effector of fenofibrate \citep{staels1998fibrates, brodie2001cetp}, and the discovery of this connection by Draphnet from bioassay and genetic data alone supports the model's biological interpretability.

\subsection{Drug Groupings Based on Shared Disease Gene Effects}

Network analysis of drug--disease gene connections identified cliques of drugs with significantly overlapping disease gene associations (Table~\ref{tab:cliques}).

\begin{table}[H]
\centering
\caption{Properties of drug cliques identified by shared disease gene associations.}
\label{tab:cliques}
\begin{tabular}{lccc}
\toprule
\textbf{Clique property} & \textbf{Mean} & \textbf{Range} & \textbf{Example} \\
\midrule
Drugs per clique & $4.2$ & $3$--$8$ & NSAIDs + fibrates \\
Shared disease genes & $6.8$ & $2$--$14$ & CETP, IL6R \\
ATC classes per clique & $2.1$ & $1$--$4$ & M01A + C10AB \\
Within-clique $\rho$ & $0.52 \pm 0.12$ & $0.31$--$0.78$ & -- \\
\bottomrule
\end{tabular}
\end{table}

These cliques frequently spanned multiple ATC therapeutic classes, revealing drug groupings based on shared downstream genomic effects rather than traditional pharmacological classification.

\section{Discussion}

Draphnet demonstrates that integrating in vitro bioassay data with genome-wide disease gene associations in a supervised framework can both predict and explain drug--phenotype associations. The model's key advantage over purely predictive approaches \citep{natarajan2014inductive, luo2017network} is interpretability: the learned interaction matrix $\mathbf{W}$ directly maps drug molecular endpoints to disease genes, enabling biological hypothesis generation.

The validation against DrugBank targets confirms that the learned projections capture real drug biology. The discovery of CETP as a shared downstream target of anti-inflammatories and PPAR$\alpha$-agonists illustrates the model's ability to identify non-obvious mechanistic connections. Drug cliques that transcend traditional pharmacological boundaries may suggest novel repurposing hypotheses and predict shared side-effect profiles \citep{pushpakom2019drug}.

Limitations include the incomplete nature of the SIDER training labels (unrecorded associations are not necessarily absent), the sparsity of ToxCast data (${\sim}$100 endpoints assayed per drug out of 1391), and the reliance on eQTL-based gene mapping in PhenomeXcan which may miss some disease genes \citep{barbeira2019exploiting}. Future work could incorporate additional data sources such as structural fingerprints or protein--protein interaction networks.

\section{Methods}

\subsection{Data Sources}
ToxCast Level 5 hit-call fractions were obtained from the EPA for 429 compounds across 1391 endpoints \citep{richard2018toxcast}. PhenomeXcan gene--phenotype z-scores were derived from S-MultiXcan \citep{barbeira2018smultixcan} applied to UK Biobank GWAS \citep{bycroft2018ukbiobank, pividori2020phenomexcan}. SIDER drug--phenotype associations \citep{kuhn2016sider} were matched to UK Biobank phenotypes via UMLS concept identifiers \citep{bodenreider2004umls}.

\subsection{Preprocessing}
ToxCast data were reduced via SoftImpute \citep{mazumder2010softimpute} with cross-validated rank and regularization parameters (5\% held-out non-missing values, MSE metric). PhenomeXcan z-scores were computed as $|z| \cdot \mathrm{sign}(z_\mathrm{consensus})$ and filtered to the top half of genes by SD.

\subsection{Model Training}
The affinity regression $\hat{\mathbf{Y}} = \mathbf{D} \cdot \mathbf{W} \cdot \mathbf{P}^\top$ was trained with L2 regularization on SIDER binary labels. Separate models were fit for side effects and indications. Cross-validation was used for hyperparameter selection.

\subsection{Validation}
Drug target significance was assessed by comparing pairwise Draphnet projection similarity of drugs sharing each DrugBank target \citep{wishart2018drugbank} against null distributions from permuted projections (rank sum test). Drug cliques were identified as fully connected subgraphs of 3+ drugs with significant overlap in disease gene associations.

\section{Conclusion}

Draphnet provides an interpretable framework for understanding how drug molecular effects propagate to disease phenotypes by learning a network connecting ToxCast bioassay endpoints to PhenomeXcan disease genes. The model predicts drug--phenotype associations, identifies downstream disease genes impacted by drug targets, and reveals biologically meaningful drug groupings that transcend traditional pharmacological classification.

\bibliographystyle{plainnat}
\bibliography{references}

\end{document}
